\documentclass[12pt]{article}
\usepackage[utf8]{inputenc}
\usepackage{amsmath,amssymb,amsthm}
\usepackage{algorithm}
\usepackage{algpseudocode}
\usepackage{geometry}
\geometry{a4paper, margin=1in}

\newtheorem{theorem}{Theorem}
\newtheorem{lemma}[theorem]{Lemma}
\newtheorem{corollary}[theorem]{Corollary}
\newtheorem{proposition}[theorem]{Proposition}
\theoremstyle{definition}
\newtheorem{definition}[theorem]{Definition}

\title{Problem 39: Integer Right Triangles}
\author{Project Euler Solution}
\date{}

\begin{document}
\maketitle

\section{Problem Statement}
For which value of $p \le 1000$ is the number of integer-sided right triangles with perimeter~$p$ maximized?

\section{Definitions}

\begin{definition}
A \emph{Pythagorean triple} $(a,b,c) \in \mathbb{Z}_{>0}^3$ satisfies $a^2 + b^2 = c^2$. It is \emph{primitive} if $\gcd(a,b,c) = 1$.
\end{definition}

\begin{definition}
For $p \in \mathbb{Z}_{>0}$, define
\[
    N(p) = \#\{(a,b,c) \in \mathbb{Z}_{>0}^3 : a \le b < c,\; a^2 + b^2 = c^2,\; a+b+c = p\}.
\]
\end{definition}

\section{Theoretical Development}

\begin{theorem}[Euclid's parametrization]
Every primitive Pythagorean triple $(a,b,c)$ with $a$ odd and $b$ even is uniquely given by $a = m^2 - n^2$, $b = 2mn$, $c = m^2 + n^2$, where $m > n > 0$, $\gcd(m,n) = 1$, and $m \not\equiv n \pmod{2}$.
\end{theorem}

\begin{lemma}[Perimeter formula]
\label{lem:perim}
For the primitive triple generated by $(m,n)$, the perimeter is $p_0 = 2m(m+n)$. The general triple scaled by $k$ has perimeter $p = 2km(m+n)$.
\end{lemma}
\begin{proof}
$a + b + c = (m^2 - n^2) + 2mn + (m^2 + n^2) = 2m^2 + 2mn = 2m(m+n)$.
\end{proof}

\begin{corollary}
$N(p) = 0$ for all odd $p$.
\end{corollary}

\begin{theorem}[Closed form for $b$]
\label{thm:b}
Given $a + b + c = p$ and $a^2 + b^2 = c^2$:
\[
    b = \frac{p(p - 2a)}{2(p - a)}.
\]
\end{theorem}
\begin{proof}
Substituting $c = p - a - b$ into $a^2 + b^2 = c^2$ and simplifying:
\begin{align*}
    a^2 + b^2 &= p^2 - 2p(a+b) + a^2 + 2ab + b^2, \\
    0 &= p^2 - 2pa - 2pb + 2ab, \\
    b &= \frac{p(p-2a)}{2(p-a)}.  \qedhere
\end{align*}
\end{proof}

\begin{lemma}[Bounds on $a$]
For $a \le b < c$ with $a + b + c = p$: $1 \le a < p/3$.
\end{lemma}
\begin{proof}
From $a \le b < c$: $p = a + b + c > a + a + a = 3a$.
\end{proof}

\begin{theorem}[Solution]
\label{thm:solution}
$\displaystyle\max_{p \le 1000} N(p) = N(840) = 8$.
\end{theorem}
\begin{proof}
By Theorem~\ref{thm:b}, $N(p)$ counts integers $a \in [1, \lfloor p/3 \rfloor)$ with $2(p-a) \mid p(p-2a)$ and $b \ge a$. The value $840 = 2^3 \cdot 3 \cdot 5 \cdot 7$ has 32 divisors, maximizing the number of valid factorizations. The 8 triples for $p = 840$ are:
\begin{center}
\begin{tabular}{ccc}
\hline
$a$ & $b$ & $c$ \\
\hline
40 & 399 & 401 \\
56 & 390 & 394 \\
105 & 360 & 375 \\
120 & 350 & 370 \\
140 & 336 & 364 \\
168 & 315 & 357 \\
210 & 280 & 350 \\
240 & 252 & 348 \\
\hline
\end{tabular}
\end{center}
Exhaustive computation over even $p \le 1000$ confirms no other perimeter achieves $N(p) \ge 8$.
\end{proof}

\section{Algorithm}

\begin{algorithmic}[1]
\Function{MaxPythagoreanPerimeter}{$P$}
    \State $\mathit{best\_p} \gets 0$, $\mathit{best\_count} \gets 0$
    \For{$p = 2$ to $P$ step $2$}
        \State $\mathit{count} \gets 0$
        \For{$a = 1$ to $\lfloor p/3 \rfloor - 1$}
            \State $b \gets p(p-2a) \;/\; 2(p-a)$
            \If{$b$ is a positive integer and $b \ge a$}
                \State $\mathit{count} \gets \mathit{count} + 1$
            \EndIf
        \EndFor
        \If{$\mathit{count} > \mathit{best\_count}$}
            \State $\mathit{best\_count} \gets \mathit{count}$, $\mathit{best\_p} \gets p$
        \EndIf
    \EndFor
    \State \Return $\mathit{best\_p}$
\EndFunction
\end{algorithmic}

\section{Complexity}

\begin{proposition}
The algorithm runs in $O(P^2)$ time and $O(1)$ space.
\end{proposition}
\begin{proof}
The outer loop runs $P/2$ times. The inner loop runs at most $P/3$ times per iteration, each costing $O(1)$. Total: $(P/2)(P/3) = P^2/6 = O(P^2)$. With $P = 1000$: $\sim 1.7 \times 10^5$ iterations.
\end{proof}

\section{Result}
\[
    \boxed{p = 840}
\]

\section{Answer}

\[
\boxed{840}
\]

\end{document}
